\section{Introduction}\label{sec:introduction}

A \emph{complex K3 surface} is a compact connected complex surface that
is simply connected and has trivial canonical bundle. Thus it carries
a nowhere-zero holomorphic two-form. Every complex K3 surface is
K\"ahler, but it need not be projective or contain any curve
\cite[Chapters~1 and~7]{Huybrechts}. This combination of a rigid
cohomological structure and potentially sparse algebraic geometry makes
K3 surfaces a natural setting in which to study holomorphic flexibility.

The \emph{convex approximation property} (CAP) for a complex manifold
asks that maps holomorphic near a compact convex subset of $\C^m$ can
be approximated there by maps defined on all of $\C^m$, for every
source dimension $m$. Forstneri\v c proved that CAP characterizes Oka
manifolds \cite{ForstnericCAP,Forstneric}. In its equivalent Stein-source
formulation, the Oka property includes deformation of continuous maps
to holomorphic maps, with approximation and holomorphic interpolation.
Our main result is the following precise approximation statement.

\begin{theorem}\label{thm:main}
Let $X$ be a complex K3 surface, with any smooth Hermitian distance
$d_X$. For every integer $m\geq1$, every nonempty compact convex set
$K\subset\C^m$, every open neighborhood $U$ of $K$, every holomorphic
map $f:U\to X$, and every $\epsilon>0$, there exists a holomorphic map
$F:\C^m\to X$ such that
\[
                  \sup_{z\in K}d_X(F(z),f(z))<\epsilon.
\]
\end{theorem}

Theorem~\ref{thm:main} answers the K3 question in
\cite[Section~8, Problem~C]{ForstnericLarussonSurvey} and proves the
positive conjecture formulated explicitly in
\cite[Introduction]{XieZhao}. Its scope includes nonprojective
surfaces and imposes no condition on the Picard rank or existence of
an elliptic fibration.

An \emph{entire curve} is a holomorphic map from $\C$, and it is an
\emph{immersion} if its derivative never
vanishes. The approximation and interpolation theory gives the
following concrete consequence, proved in Section~\ref{sec:dense-entire}.

\begin{corollary}\label{cor:dense}
Let $S$ be a complex K3 surface. For every $x\in S$
and every nonzero vector $v\in T_xS$, there is a holomorphic immersion
$f:\C\to S$ such that
\[
 f(0)=x,\qquad f'(0)=v,\qquad \overline{f(\C)}=S,
\]
where closure is taken in the ordinary complex topology. If $S$ is
projective, then $f(\C)$ is Zariski dense.
\end{corollary}

Here the tangent vector, not just its line, is prescribed. For projective
$S$, the density conclusion confirms the K3 case of Campana's
prediction of Zariski-dense entire curves on special projective
manifolds \cite{Campana}. Here specialness means the absence of a
meromorphic fibration whose positive-dimensional base, endowed with
the orbifold structure induced by the multiple fibers, is of general
type.
See \cite[Conjecture~1.3]{CCR} for the prediction and
\cite[Proposition~8.9]{CCR} for the specialness of projective varieties
with trivial canonical bundle. The ordinary density and prescribed
first jet are additional conclusions.

Section~\ref{sec:surface-consequences} derives the Oka property for
Enriques surfaces and hence for all minimal compact complex surfaces
of Kodaira dimension zero. It also obtains global dominating sprays
on projective K3 and Enriques surfaces. A separate class-VII
classification there combines earlier Oka results with the global
spherical shell theorem of the companion manuscript \cite{OpenAIGSS};
that theorem is not an input to the K3 argument.

\subsection{Earlier results and the remaining obstacles}

The use of holomorphic sprays and approximation to realize topological
flexibility is central to Gromov's Oka principle for elliptic bundles
\cite{GromovOka}. Forstneri\v c's convex approximation characterization
\cite[Theorem~0.1]{ForstnericCAP} makes CAP an effective criterion for
this flexibility, while his interpolation theorem
\cite[Corollary~1.3]{ForstnericInterpolation} links it to prescribed
holomorphic data on subvarieties. Our proof uses local chart sprays
and splitting arguments in this tradition, with the uniform estimates
needed for the strip and convex-face constructions made explicit.

Holomorphic flexibility on K3 surfaces has been studied through entire
curves, domination by affine space, and approximation. Buzzard and Lu
proved that elliptic and Kummer K3 surfaces are holomorphically
dominable by $\C^2$ \cite[Propositions~4.4--4.5]{BuzzardLu}.
Dominability means the existence of a map $\C^2\to X$ whose
differential is surjective somewhere. Forstneri\v c and L\'arusson
strengthened this for every Kummer surface to \emph{strong
dominability}: such a map can pass through each prescribed target
point with surjective differential
\cite[Corollary~3]{ForstnericLarussonSurfaces}. These properties
concern individual maps; CAP asks for approximation of every local
map in every source dimension.

For arbitrary complex K3 surfaces, Kamenova, Lu and Verbitsky proved
that the Kobayashi pseudodistance vanishes
\cite[Corollary~2.2]{KamenovaLuVerbitsky}; the period-dynamics argument
is corrected in \cite[Section~4]{VerbitskyErratum}. This removes an
intrinsic obstruction to holomorphic flexibility, but does not
provide approximation or interpolation.

Alarc\'on and Forstneri\v c introduced the Oka-1 property and established
it for Kummer surfaces and elliptic K3 surfaces
\cite[Proposition~8.4 and Corollary~8.6]{AlarconForstneric}.
Oka-1 concerns approximation and discrete jet interpolation from open
Riemann surfaces, within prescribed continuous homotopy classes. Its
full formulation is recalled below.
The source dimension distinction matters: the one-dimensional
property alone does not give Theorem~\ref{thm:main}.

For projective K3 geometry, Chen, Gounelas and Liedtke developed
existence results for curves of prescribed geometric genus
\cite{ChenGounelasLiedtke}. Chen and Gounelas then proved that every
projective K3 surface contains genus-one curves of unbounded
self-intersection whose smooth normalizations vary in moduli
\cite[Theorem~A]{ChenGounelas}. We use two of their families with
different ample degrees. Their complete genus-one fibers supply
holomorphic flows for all complex times, and the difference in degrees
forces their tangent directions to be generically independent.

More recently, Xie and Zhao proved that smooth $(2,2,2)$ hypersurfaces
in $(\PP^1)^3$ are Oka and established dense $G_\delta$ sets of Oka
K3 periods and corresponding local deformation parameters
\cite[Theorems~A, B, and~E]{XieZhao}. Their work combines geometric
constructions of flexibility with arithmetic dynamics of periods and
Torelli theory. We use the same period-theoretic framework, including
the corrected orbit alternatives of Verbitsky
\cite{VerbitskyErratum}, with different geometric constructions.
A dense collection of Oka surfaces does not by itself settle the
property for each surface: the transfer must address every possible
rational subspace of its period plane.

Two analytic obstacles arise even before that transfer. Local
approximation in one variable must retain arbitrary holomorphic
parameters and then imply approximation on arbitrary convex sets.
Moreover, the geometric constructions produce long chains of annuli;
errors at their joins must be corrected without exhausting a fixed
positive transverse width. We handle the first issue by a polynomial
change of source coordinates and gluing estimates measured on real
faces. We handle the second by exact symplectic sewing with bounds
independent of the chain lengths.

\subsection{The proof and its organization}

The analytic and geometric parts meet at a local operation, denoted
$L$. It enlarges a planar disc across an attached disc, for maps with
image in a sufficiently small target neighborhood, while retaining
any prescribed smaller polydisc of passive holomorphic parameters.
Section~\ref{sec:local} defines $L$ precisely and obtains it from two
strip maps, one defined on $\C\times\Delta$ and the other on
$\Delta\times\C$, with transverse complete directions. Aligning
their coordinate axes exactly makes the overlap error arbitrarily
small. A bounded additive splitting and a contraction correct it.
The same section propagates $L$ across the center of a small disc
whose boundary already lies in the locus where $L$ holds.

Sections~\ref{sec:gluing} and~\ref{sec:globalization} prove that $L$
at every point implies CAP in every dimension
(Theorem~\ref{ana:cap}). A polynomial automorphism concentrates the
part of an enlarged polydisc outside the original domain into one
long coordinate and bounded transverse coordinates. The planar
operation with passive parameters then gives approximation on
polydiscs. To pass to an arbitrary convex set, we remove the defining
real faces of a containing polytope one at a time. The necessary
gluing estimate uses H\"older accuracy on each real face, without
requiring a fixed complex thickness on which that accuracy holds.

Section~\ref{sec:projective} constructs the two complete directions
on projective K3 surfaces using the Chen--Gounelas families. Propagation
across their algebraic exceptional set gives $L$ everywhere. This
handles one case of the final period argument and also exhibits a
direct source of the required local geometry.

The nonprojective cases require strips on open regions of period
space. Section~\ref{sew:section} first establishes the uniform sewing
theorem. On an annulus, a closed holomorphic one-form can have a
nonzero period. Exactness of the symplectic transition removes this
obstruction: the transverse constant Fourier coefficient becomes
quadratic in the error. The other Fourier modes admit a splitting
whose bounds are independent of the number and lengths of the
annuli. A quadratic iteration then retains a positive transverse
radius along an infinite chain.

To describe the geometric transfer, fix the K3 lattice $\Lambda$,
the integral second-cohomology lattice with its intersection form.
A \emph{marking} identifies $H^2(X,\Z)$ with $\Lambda$. The real
and imaginary parts of a nonzero holomorphic two-form span an
oriented positive two-plane, the \emph{period} of the marked surface.
Its rational directions are the vectors in its intersection with
$\Lambda\otimes\Q$.

Section~\ref{deg:section} constructs a nonempty open region in the
full period domain on which every surface has $L$ everywhere. It
smooths a union of two quadrics and sews annuli along their common
elliptic curve. The chart estimates persist under unrestricted small
K3 deformations. Section~\ref{rat:section} constructs a second kind
of region: for each fixed positive integral vector $v$, a relatively
open set among the period planes containing $v$. Here an isotrivial
genus-one fibration supplies the central model. Conservation of
imaginary symplectic action keeps all recentered annuli away from
singular fibers. Two independent cycles supply distinct strip leaves
after deformation and hence the transverse directions needed for $L$.

Finally, Section~\ref{dyn:section} applies Ratner's orbit theorem and
Verbitsky's corrected rational-direction analysis. Periods with no
rational direction have orbits meeting the first open region. Periods
with exactly one rational direction have orbits, under its stabilizer,
meeting the corresponding second region. Rational period planes give
projective K3 surfaces, already treated above. Torelli theory
transfers $L$ along these orbits, and Theorem~\ref{ana:cap} finishes
the proof. Section~\ref{sec:dense-entire} derives interpolation,
Corollary~\ref{cor:dense}, and the further surface consequences.

\paragraph{Conventions and standard inputs.}
Maps on a compact set are always defined on a neighborhood of that set.
Approximation may lose any prescribed positive amount of an auxiliary
domain margin. A strip means $\C\times\Delta_b$, with no injectivity
requirement unless stated; local regularity is specified at each use.
The analytic background consists of Stein theory, scalar approximation
and the $L^2$ $\bar\partial$ theorem
\cite{Forstneric,Hormander,Siu}. We recall the needed K3 period,
Torelli, and nef-pencil results at their applications
\cite[Chapters~2, 6--8, and~11]{Huybrechts}. The arithmetic inputs
are likewise stated where used. The local approximation, sewing,
geometric constructions, and reductions between them are proved here.
